%============================================================
\chapter{Sự Tha Thứ (The Power of Forgiveness)}
\begin{center}
  \textit{\color{truyenblue}Forgiveness is not an occasional act;\\
  it is a constant attitude.}\\
  \textit{(Tha thứ không phải là hành động đôi khi;\\
  đó là thái độ thường trực.)}\\[4pt]
  \small\color{truyengold}--- Martin Luther King Jr. ---
\end{center}
\ngancach

\section{Nelson Mandela Ra Tù Không Oán Thù}
\begin{truyen}{Nelson Mandela -- 27 Năm Tù Và Trái Tim Không Hận Thù}{Nam Phi, Thế Kỷ XX}
\chuhoa{N}{}N elson Mandela bị giam 27 năm trong nhà tù Robben Island vì đấu tranh chống phân biệt chủng tộc.\\
\textit{(Nelson Mandela was imprisoned for 27 years on Robben Island for fighting against apartheid.)}

Khi được thả, người ta hỏi: ``Ông có căm thù những người giam cầm ông không?''\\
\textit{(When released, people asked: ``Do you hate the people who imprisoned you?'')}

Mandela lắc đầu: ``Nếu tôi còn căm thù khi bước ra cánh cổng này, tôi vẫn còn là tù nhân trong lòng.''\\
\textit{(Mandela shook his head: ``If I still harbored hatred as I walked out these gates, I would still be a prisoner within myself.'')}

Ông trở thành Tổng thống Nam Phi và mời những cai tù cũ ngồi hàng đầu trong lễ nhậm chức.\\
\textit{(He became President of South Africa and invited his former prison guards to sit in the front row at his inauguration.)}

Ông nói: ``Tha thứ không phải là quà tặng cho kẻ đã làm hại ta -- đó là quà tặng ta trao cho chính mình.''\\
\textit{(He said: ``Forgiveness is not a gift for those who have wronged us -- it is a gift we give to ourselves.'')}

\end{truyen}
\begin{baihoc}
  Tha thứ không phải là bào chữa cho kẻ xấu. Tha thứ là giải phóng bản thân khỏi gánh nặng oán thù.\\
  \textit{(Forgiveness is not an excuse for wrongdoers. Forgiveness is liberating yourself from the burden of resentment.)}
\end{baihoc}

\section{Phật Thích Ca Và Bài Học Mũi Tên}
\begin{truyen}{Mũi Tên Thứ Hai}{Phật Giáo}
\chuhoa{Đ}{}Đ ức Phật dạy về nỗi đau bằng ẩn dụ hai mũi tên bắn vào người.\\
\textit{(The Buddha taught about pain using the metaphor of two arrows shot at a person.)}

Mũi tên thứ nhất là sự kiện thực tế xảy đến: mất mát, bị xúc phạm, bị tổn thương.\\
\textit{(The first arrow is the actual event that happened: loss, insult, injury.)}

Mũi tên thứ hai -- đau hơn nhiều -- là phản ứng của ta: tức giận, oán thù, tiếc nuối mãi mãi.\\
\textit{(The second arrow -- far more painful -- is our reaction: anger, resentment, endless regret.)}

Người thường bắn vào mình mũi tên thứ hai tự nguyện và còn giữ nó ở đó lâu hơn.\\
\textit{(Ordinary people shoot the second arrow into themselves voluntarily and keep it there even longer.)}

Đức Phật dạy: ``Trong cuộc đời, mũi tên thứ nhất không thể tránh, nhưng mũi tên thứ hai là lựa chọn của ta.''\\
\textit{(The Buddha taught: ``In life, the first arrow cannot always be avoided, but the second arrow is our own choice.'')}

\end{truyen}
\begin{baihoc}
  Tha thứ là rút mũi tên thứ hai ra khỏi trái tim mình. Đau khổ từ việc không tha thứ luôn tự làm ta tổn thương hơn.\\
  \textit{(Forgiveness is pulling the second arrow out of your own heart. The suffering from not forgiving always hurts us more.)}
\end{baihoc}

\section{Câu Chuyện Hai Anh Em Người Đức}
\begin{truyen}{Hai Anh Em Và Mối Thù Ba Mươi Năm}{Câu Chuyện Hiện Đại}
\chuhoa{H}{}H ai anh em người Đức cùng kinh doanh giày dép, sau một tranh cãi nhỏ đã chia đôi cơ sở.\\
\textit{(Two German brothers who ran a shoe business together, after a small quarrel, split the company in two.)}

Adolf lập Adidas; Rudolf lập Puma -- hai thương hiệu đối địch chỉ cách nhau một dòng sông.\\
\textit{(Adolf founded Adidas; Rudolf founded Puma -- two rival brands separated by only a river.)}

Hai anh em không nói chuyện nhau suốt 30 năm -- bạn bè, gia đình bị buộc chọn phe.\\
\textit{(The two brothers did not speak for 30 years -- friends and family were forced to choose sides.)}

Cả thị trấn Herzogenaurach bị chia đôi bởi mối thù không có hồi kết.\\
\textit{(The entire town of Herzogenaurach was divided in two by the endless feud.)}

Trước khi mất, Rudolf được hỏi điều hối tiếc nhất: ``Tôi đã lãng phí 30 năm tốt nhất của mình vào mối thù không đáng.''\\
\textit{(Before he died, Rudolf was asked what he regretted most: ``I wasted my 30 best years on a feud not worth having.'')}

\end{truyen}
\begin{baihoc}
  Mỗi ngày không tha thứ là một ngày ta trao sức mạnh của mình cho người đã làm ta tổn thương. Tha thứ là lấy lại sức mạnh đó.\\
  \textit{(Every day of not forgiving is a day we give our power to the person who hurt us. Forgiveness is taking that power back.)}
\end{baihoc}

\section{Vua Lý Thái Tông Và Người Trộm Lúa}
\begin{truyen}{Lý Thái Tông Xử Án Nhân Từ}{Việt Nam, Thế Kỷ XI}
\chuhoa{M}{}M ột lần vua Lý Thái Tông đang đọc sách thì nghe tin có tên trộm lúa trong kho hoàng gia bị bắt.\\
\textit{(Once, King Ly Thai Tong was reading when news came that a grain thief had been caught in the royal storehouse.)}

Quan tâu tên trộm phải bị chặt tay theo lệ, vua hỏi: ``Hắn ăn trộm vì nghèo hay vì tham?''\\
\textit{(Officials reported the thief should have his hand cut off by custom; the king asked: ``Did he steal out of poverty or greed?'')}

Điều tra ra hắn có vợ ốm nặng, con nhỏ đói khóc, đã ba ngày không có gì ăn.\\
\textit{(Investigation revealed he had a gravely ill wife, hungry crying children, and had not eaten for three days.)}

Vua phán: ``Kẻ trộm vì đói là lỗi của triều đình -- ta chưa lo được cho dân. Tha cho hắn.''\\
\textit{(The king declared: ``A thief driven by hunger is the fault of the court -- we have not cared for the people. Pardon him.'')}

Rồi vua cấp thêm lúa và thuốc cho gia đình hắn, ra chiếu mở kho cứu đói cả vùng.\\
\textit{(Then the king provided grain and medicine for the man's family and issued an edict to open the granaries for famine relief throughout the region.)}

\end{truyen}
\begin{baihoc}
  Tha thứ đích thực đi cùng với hành động tích cực. Người hiền minh không chỉ tha tội mà còn giải quyết nguyên nhân gốc rễ.\\
  \textit{(True forgiveness comes with positive action. The wise do not only pardon the sin but also address its root cause.)}
\end{baihoc}

\section{Corrie Ten Boom Và Tên Lính Đức}
\begin{truyen}{Corrie Ten Boom -- Tha Thứ Điều Không Thể Tha Thứ}{Hà Lan -- Đức, Thế Kỷ XX}
\chuhoa{C}{}C orrie Ten Boom sống sót qua trại tập trung Đức Quốc Xã -- chứng kiến chị gái chết trong tù.\\
\textit{(Corrie Ten Boom survived the Nazi concentration camps -- witnessing her sister die in captivity.)}

Sau chiến tranh, một hôm bà gặp lại tên lính đã canh gác phòng tắm, đã xem bà và chị gái như đồ vật.\\
\textit{(After the war, one day she encountered the very guard who had watched over the shower room, who had treated her and her sister like objects.)}

Tên lính giơ tay bắt tay: ``Cảm ơn bà vì thông điệp về sự tha thứ -- Chúa đã tha tội cho tôi chưa?''\\
\textit{(The guard extended his hand: ``Thank you for your message of forgiveness -- has God forgiven me?'')}

Bà đứng đó, tim như đá lạnh. Bà không thể đưa tay ra. Bà cầu nguyện thầm.\\
\textit{(She stood there, heart like cold stone. She could not extend her hand. She prayed silently.)}

Rồi bà đưa tay ra -- không phải vì cảm xúc, mà vì lựa chọn tha thứ -- và ngay lúc đó sự ấm áp tràn vào.\\
\textit{(Then she extended her hand -- not from emotion but from a choice to forgive -- and in that moment warmth flooded in.)}

Bà viết: ``Tha thứ là hành động ý chí, không phải cảm xúc. Hành động đến trước, cảm xúc theo sau.''\\
\textit{(She wrote: ``Forgiveness is an act of will, not emotion. The action comes first, feelings follow.'')}

\end{truyen}
\begin{baihoc}
  Tha thứ không phải cảm giác -- đó là quyết định. Khi ta quyết định tha thứ, trái tim sẽ từ từ lành lại.\\
  \textit{(Forgiveness is not a feeling -- it is a decision. When we decide to forgive, the heart slowly heals itself.)}
\end{baihoc}
